\section{Introduction}
\label{sec:introduction}

The rational-distance problem for a square asks for a point whose distances
from all four vertices are rational. We consider the strict interior of the
unit square and retain the four distances as part of the arithmetic data.
Writing the vertices in the order $00,10,01,11$, the equations are
\[
 \begin{aligned}
 r_0^2&=x^2+y^2,&r_1^2&=(1-x)^2+y^2,\\
 r_2^2&=x^2+(1-y)^2,&r_3^2&=(1-x)^2+(1-y)^2,
 \end{aligned}
 \qquad 0<x,y<1,
\]
with $r_i\in\Q_{>0}$. Differences of the squared distances recover $x$
and $y$, so this is a rational-point problem with a specified real chamber.

The geometry of rational-distance varieties has been studied from several
directions. Berry relates the square to a Kummer surface and a branched
double cover, and studies associated curves of genus five
\cite{berry1990}; his work on triangles develops the corresponding
three-distance geometry \cite{berry1992}. Bremner and Ulas give
constructions for rational-distance problems attached to several geometric
objects \cite{bremnerulas2016}. Corvaja, Turchet and Zannier study distance
varieties for general configurations of rational points, including their
surfaces and product-quotient descriptions
\cite{corvajaturchetzannier2025}. The present paper develops the arithmetic
of explicit covering curves, the squareclasses carried by one ordered
source, and congruences that hold at every denominator height.

The first part concerns the primitive integral equations. Let $L$ be the
least common denominator of $x,y$, and write $A_i=Lr_i$. We prove that
$24\mid L$, that the common radius content $G=\gcd(A_0,A_1,A_2,A_3)$
satisfies $G^2\mid L$, and that the six reduced denominators of the pairwise
radius sums and differences satisfy exact gcd and product identities
(Theorem~\ref{thm:int:pair-denominators}). Signed radius identities compare
each prime-power divisor of $L$ with the size of the square. They imply,
in particular, that $L$ has a prime divisor other than $2$ or $3$
(Corollary~\ref{cor:int:six-smooth}).

Gaussian factorization assigns four canonical, pairwise coprime
squarefree integers $m_0,\ldots,m_3$ to the corners. Their product is
the squarefree part of
\[
 (r_0r_3+r_1r_2)^2-1.
\]
All primes in these integers are $1$ or $7$ modulo $8$. A reciprocity
identity for each pair of corners retains the common denominator and the
sign of the Gaussian norm. For the two opposite pairs it gives
\[
 m_0m_3\equiv1\text{ or }15\pmod{16},\qquad
 m_1m_2\equiv1\text{ or }15\pmod{16}
\]
(Theorem~\ref{thm:profile-main}). We also classify the residue pattern
arising when $v_p(L)=2v_p(G)$. An elliptic residue curve identifies the
exceptional primes $5$ and $13$; imposing the additional product-square
condition gives the refined exceptional set $\{5,13,17\}$
(Theorems~\ref{thm:crit-basic} and~\ref{thm:crit-refined}).

The second part gives explicit rational correspondences. An ordered pair
$Q=(a,b)$, $R=(c,e)$ on
\[
 E_n:\quad Y^2=X(X-n)(X+n)
\]
encodes the original point through two ordered product squareclasses,
the two ordinates, and the single further square
\[
 \frac{a+c}{2(c+n)}.
\]
Theorem~\ref{thm:ell-exact-source} gives the forward and inverse maps,
including all signs and strict inequalities. Fixing a signed sum or
difference of two distances instead gives a smooth curve of genus five.
The resulting fixed-pair finiteness theorem follows from the explicit
branch calculation and Faltings' theorem
(Theorem~\ref{thm:ell-fixed-pair-finite}).

The product parameter $t=r_0r_3+r_1r_2$ leads to a genus-nine curve $D_t$
with four ordered square roots over a common coordinate $\ell$. Its
valuation conditions constrain one global class $[\ell]$. A generally
nonsymmetric Legendre-symbol matrix then reduces the possible classes to
two, and a local identity at $37$ excludes both under an explicit residue
condition. This proves an infinite family of all-height no-point results
(Theorem~\ref{thm:common-squareclass} and
Corollary~\ref{cor:common-infinite}). The associated unramified double
covers have genus seventeen. Nine explicit elliptic quotients describe
$\Jac(D_t)$ and yield rank constraints on parameters admitting rational
points (Theorems~\ref{thm:nine-isogeny}
and~\ref{thm:nine-rank-strata}).

Further descent must compare local data arising from the same global
source. We express this requirement through a character map and a
secondary obstruction in its cokernel. Keeping the complete compensation
coset gives the criterion of Theorem~\ref{thm:cs-matrix}. A short
arithmetic certificate applies it, together with a higher quotient, to
exclude the entire kernel $n=15569$
(Corollary~\ref{cor:cs-15569-layer}). A separate construction attaches
a proper quaternion class to a complete genus-five source. Its local
evaluations give an explicit parity criterion through Hilbert reciprocity
(Theorem~\ref{thm:hp-uniform}). An explicit conic parametrization yields
an infinite family of complete genus-five failures of the Hasse principle
(Theorem~\ref{thm:hp-infinite}).

The final section treats the fixed opposite-distance difference $1/5$
in detail. Its genus-five source has an unramified genus-three quotient
which is a smooth plane quartic. We determine all twenty-eight
bitangents, their contact divisors and Galois action, obtaining the
odd-theta algebra
\[
 \Q^8\times\Q(\sqrt2)^4\times\Q(\sqrt{51})^4
 \times\Q(\sqrt2,\sqrt{51}).
\]
The calculation supplies explicit theta fractions of the type used in
generalized descent \cite{bruinpoonenstoll2016,misme2026}. A polynomial
identity identifies the original double cover with a specified difference
of theta characteristics. We then prove a local-solubility theorem for
every finite Abel--Jacobi multiplication cover through a rational boundary
point. Its application preserves the original double cover and a strict
real point (Corollary~\ref{cor:theta-surviving-covers}).

For the models encoding the original four-distance configurations, we give
explicit reconstruction formulas and the inequalities defining the strict
real chamber. Finite residue calculations
are given as explicit identities, congruences or complete small tables.
We use the classical torsion theorem of Mazur, the finiteness theorem of
Faltings, and standard facts about elliptic curves, covers and Brauer groups
in the statements where they enter. The appendix describes the compact
four-radius surface and the Brauer evaluations associated with its real
components.

\paragraph{Companion material.}
The paper sources and a partial Lean formalization are archived at
\url{https://doi.org/10.5281/zenodo.23638745}.
The archive identifies the verified statements and the pinned build environment.
